\VerdierPhysicalPageAnchor{0043}%
\VerdierSemanticAnchor{ch1-1-3-8-composition-cont}
de la composition dans la catégorie graduée $\mathcal C[G]$ associée à
la $G^\circ$-catégorie $(\mathcal C,G^\circ,T,c)$ (1.3.2.5).

\VerdierSemanticAnchor{ch1-1-3-9}
\noindent\textbf{1.3.9.}\quad
Soit $(\mathcal C,G,T,c)$ une $G$-catégorie. La catégorie opposée à
$(\mathcal C,G,T,c)$ est une $G^\circ$-catégorie ($G^\circ$ groupe
opposé à $G$) $(\mathcal C^\circ,G^\circ,T^\circ,c^\circ)$ définie de
la manière suivante. La catégorie $\mathcal C^\circ$ est la catégorie
opposée à $\mathcal C$. Les foncteurs $T^\circ(g)$, $g\in G^\circ$,
sont définis par :

\VerdierSemanticAnchor{ch1-1-3-9-eq1}
\begin{equation}
  \begin{aligned}
    T^\circ(g)X &= T(g^{-1})X,
      & X&\in\operatorname{Ob}(\mathcal C^\circ)
             =\operatorname{Ob}(\mathcal C),\\
    T^\circ(g)u &= T(g^{-1})u,
      & u&\in\operatorname{Fl}(\mathcal C^\circ)
             =\operatorname{Fl}(\mathcal C),
  \end{aligned}
  \tag{1.3.9.1}
\end{equation}
et les isomorphismes $c^\circ(g,h)$ de $\mathcal C^\circ$ sont définis
par :

\VerdierSemanticAnchor{ch1-1-3-9-eq2}
\begin{equation}
  \begin{aligned}
    c^\circ(g,h)&\colon
      T^\circ(h\perp g)\longrightarrow T^\circ(h)T^\circ(g),\\
    c^\circ(g,h)&=c^{-1}(g^{-1},h^{-1}).
  \end{aligned}
  \tag{1.3.9.2}
\end{equation}
(on désigne par $\perp$ la composition dans $G^\circ$).

\VerdierSemanticAnchor{ch1-1-3-10}
\noindent\textbf{Remarque 1.3.10.}\quad
Soit $\mathcal C[G^\circ]\longrightarrow G^\circ$ la catégorie graduée
de type $G^\circ$ associée à la catégorie $(\mathcal C,G,T,c)$. Cette
catégorie est munie d'un clivage canonique. Soit alors
$\mathcal C[G^\circ]^\circ\longrightarrow G$ la catégorie cofibrée
obtenue en passant aux catégories opposées. Cette catégorie est munie
d'un coclivage canonique et, comme $G$ est un groupe, cette catégorie
est aussi fibrée sur $G$. On en déduit un clivage canonique sur
$\mathcal C[G^\circ]^\circ\longrightarrow G$ et par suite une
$G^\circ$-catégorie qui est canoniquement isomorphe à la catégorie
opposée à la catégorie $(\mathcal C,G,T,c)$, ainsi qu'on le vérifie
immédiatement.

\VerdierSemanticAnchor{ch1-1-3-11}
\noindent\textbf{1.3.11.}\quad
Soient $(\mathcal C',G',T',c')$ une $G'$-catégorie et
$(\mathcal C'',G'',T'',c'')$ une $G''$-catégorie. La catégorie
$\mathcal C'\times\mathcal C''$ est munie canoniquement d'une structure
de $G'\times G''$-catégorie de la manière évidente.

\VerdierSemanticAnchor{ch1-1-3-11-product}
La $G'\times G''$-catégorie obtenue sera appelée le \emph{produit} de
la $G'$-catégorie $(\mathcal C',G',T',c')$ par la $G''$-catégorie
$(\mathcal C'',G'',T'',c'')$.

\VerdierSemanticAnchor{ch1-1-3-12}
\noindent\textbf{1.3.12.}\quad
Soit $(\mathcal C,G^\circ,T,c)$ une $G^\circ$-catégorie. On utilisera
le plus souvent dans la pratique une notation à droite pour désigner
les opérations de $G^\circ$ sur $\mathcal C$. On posera pour tout objet
$X$ de $\mathcal C$ :

\VerdierSemanticAnchor{ch1-1-3-12-eq1}
\begin{equation}
  T(g)X=X[g]
  \tag{1.3.12.1}
\end{equation}
et de même pour tout morphisme $u$ de $\mathcal C$ :

\VerdierSemanticAnchor{ch1-1-3-12-eq2}
\begin{equation}
  T(g)u=u[g].
  \tag{1.3.12.2}
\end{equation}
Les isomorphismes $c(g,h)$ définissent donc, pour tout objet $X$, des
isomorphismes :

\VerdierSemanticAnchor{ch1-1-3-12-eq3}
\begin{equation}
  c(g,h)(X)\colon X[gh]\longrightarrow X[g][h].
  \tag{1.3.12.3}
\end{equation}

\VerdierSemanticAnchor{ch1-1-3-12-notation}
Cette notation a en particulier l'avantage de transformer les
opérations de $G^\circ$ sur $\mathcal C$ en des opérations à droite de
$G$ sur $\mathcal C$.
